\chapter{Conclusion and Future Work}
\label{sec:conclusion}

\section{Conclusion}
SiliconScribe works as an execution-grounded RTL assistant rather than a scripted
demonstration. A plain-English prompt is parsed into a specification, which the user
can review and edit; RTL and an independent, specification-derived testbench are
generated; Icarus Verilog executes them; and failures produce a structured,
line-level diagnosis that is patched and re-simulated. The simulator decides PASS or
FAIL, never the model. Around that loop the system measures how much a PASS is worth:
testbench strength by mutation testing, measured toggle coverage, a second-opinion
testbench after a testbench-changing repair, an optional separate model for the
testbench, and a pre-layout timing and area estimate on the SkyWater 130\,nm library.

In a live evaluation with one free model, 14 of 15 prompts completed: 8 passed on the
first attempt and 6 after one self-repair, and hand-written reference testbenches
confirmed that all 14 final designs are correct. The fifteenth prompt timed out three
times and is reported as no result. The tool itself is checked by 436 backend and 211
frontend tests, a 27-of-27 mutation sweep, a gate-level equivalence check of the
timing stage on 19 designs, and an accessibility audit with no violations.

Because the pipeline is model-agnostic, its quality is expected to rise with better models (a larger model already turned an incorrect priority encoder into a first-attempt PASS), and the remaining work is mostly infrastructure and stronger verification rather than redesign (Section~\ref{sec:scale}).

The most important engineering lesson is that \emph{honesty is a feature}. The
system labels compiler repairs, refuses to delete checks, refuses to swap designs
mid-run, refuses to time circuits it cannot time honestly, reports a stalled model as
``no result'', reports when two models disagree, and distinguishes a scripted offline
fix from a live diagnosis. Those choices make the PASS meaningful. A PASS is still
evidence, not proof.

\section{Limitations}
\label{sec:limits}

\begin{enumerate}[label=\textbf{L\arabic*.},ref=L\arabic*]
\item \textbf{Residual common-mode model risk.} The testbench never sees the RTL,
but unless a separate testbench model is chosen, the same model interprets the
specification for both. If it misreads the specification identically in both, the
pair can still agree while being wrong. The specification is also written by a model
unless the user reviews it. The separate testbench model, the specification review
and the second opinion reduce this risk; they do not remove it, and the application
does not decide automatically which side of a disagreement is wrong.
\item \textbf{Strength is a sample.} The strength score uses up to 12 syntactic
one-token mutants. A survivor may be an equivalent mutant, and a high score does not
prove the testbench complete.
\item \textbf{Only toggle coverage is measured.} Line and branch coverage are not
measured; Verilator could not build most of the generated testbenches. Any other
coverage figure is what the testbench reports about itself and is labelled as such.
\item \textbf{No assertions or formal checks.} There are no SVA properties and no
formal equivalence check; the only oracle is the simulated testbench.
\item \textbf{Timing is a pre-layout estimate.} One typical corner, no placement or
routing (no wire delay), an ideal clock, zero I/O delays and setup analysis only.
Power is not reported. Designs with tri-states, latches, loops or derived clocks are
refused rather than timed. It is not sign-off.
\item \textbf{Small evaluation.} The evaluation used small blocks (under about 100
lines), one run per prompt and mainly one model, so it does not support pass-rate
comparisons with published benchmarks.
\item \textbf{Free models are slow and quota-limited.} A live run took a median of
360~s, some free-queue models took minutes per reply, Cloudflare's free allowance is
10,000 neurons a day, and models are retired without notice. GitHub Copilot through
OpenCode is available only as \code{github-copilot/auto}, which cannot name the model
it used.
\item \textbf{UVM export is not compiled.} The 16-file UVM 1.2 bundle targets a
commercial simulator; Icarus cannot compile UVM and no commercial licence was
available. It is a structural export, not a verified UVM testbench.
\item \textbf{Not a hardened sandbox.} Generated Verilog runs as a local \code{vvp} subprocess with the server's privileges. The sandbox policy (Section~\ref{sec:sandbox}), the scrubbed environment, the resource limits and the output cap reduce the attack surface (the memory limit is enforced on Linux only), but a static scan of Verilog cannot be airtight; multi-user deployment requires a container or microVM per simulation.
\item \textbf{Single-node runtime.} The rate limiter is in-process, persistence is a local SQLite file and simulations share one host; concurrency is bounded by the thread pool rather than by a queue (Section~\ref{sec:scale}).
\end{enumerate}

\section{Future Work}
\begin{enumerate}[label=\textbf{\arabic*.}]
\item \textbf{Multi-model comparison} on the fifteen evaluation prompts, with the
reference check and the strength score reported per model.
\item \textbf{Assertions and formal checks:} SVA properties as a second oracle
alongside the self-checking testbench, and formal equivalence where a reference
exists.
\item \textbf{Line and branch coverage} with a simulator that supports the generated
testbenches.
\item \textbf{Place and route with OpenROAD} for post-layout timing with real wire
delay.
\item \textbf{Gate-level visualisation} of the mapped netlist in the schematic tab.
\item \textbf{Production runtime:} container isolation per simulation and a job queue
for multi-user use, with PostgreSQL and authentication (Section~\ref{sec:scale}).
\item \textbf{Code structure:} split the pipeline function (about 500 lines) into
smaller stages that can be tested on their own.
\end{enumerate}
